
Section~\ref{sec:inference} states what each of the reporting objects is and
why it is that and not something else. This appendix carries the constructions
themselves, so that the section can be read for the argument and this one for
the arithmetic. What is here is the detail a reader would need to reimplement
the objects without reading the code.

\subsection*{The multiplier bootstrap for the max-$t$ critical value}
\label{app:multiplier}

Let $\mathcal{F}$ be the declared family, $V^{(b)}_c$ the $b$th bootstrap draw
at cell $c$, $\bar V_c$ the draw mean and $\hat\sigma_c$ the per-cell
bootstrap standard error. Standardise the draw matrix,
\[
Z_{b,c} \;=\; \frac{V^{(b)}_c - \bar V_c}{\hat\sigma_c},
\qquad b = 1,\dots,B, \quad c \in \mathcal{F},
\]
and let $e_1,\dots,e_B$ be independent standard normals drawn afresh for each
replicate. The \emph{multiplier} statistic is
\[
T^{(\mathrm{mult})} \;=\; \max_{c \in \mathcal{F}}
\Bigl| B^{-1/2}\textstyle\sum_{b} e_b Z_{b,c} \Bigr| .
\]
Conditionally on the draws, $B^{-1/2}\sum_b e_b Z_{b,\cdot}$ is a centred
Gaussian vector whose covariance is exactly the empirical covariance of the
standardised draws, so the distribution of $T^{(\mathrm{mult})}$ approximates
that of the max-$t$ statistic without any further refitting
\citep{chernozhukov2013gaussian}. The critical value $q_{1-\alpha}$ is the
$(1-\alpha)$ quantile of \nMultiplier\ such replicates, generated in chunks so
that the ${\nMultiplier}\times B$ normal matrix is never held at once.

Two properties are worth stating because they are the reason the construction
is used here. The precision of $q$ is governed by the number of multiplier
replicates and not by $B$, so a design whose fitting budget is fixed by the
cost of refitting an entire pipeline can still resolve its critical value; and
because the multiplier draws are conditional on the observed $Z$, the
construction inherits whatever dependence the moving-block resampling induced
across cells, which is what a family-wise statement over a correlated surface
requires. The approximation is asymptotic in the same sense the max-$t$ band
is: it needs the high-dimensional central limit theorem to hold over a family
whose size is comparable to the draw count, which Section~\ref{sec:lim} states
as an assumption rather than a guarantee.

The Monte Carlo interval for $q$ itself is retained. Writing
$[q^-, q^+]$ for the order-statistic interval of the multiplier quantile, a
cell is called \emph{resolved} only if its band excludes zero at $q^+$, so
residual Monte Carlo noise leaves cells unresolved rather than resolving them
in the wrong direction.

\subsection*{Fieller's construction for the reduction}
\label{app:fieller}

For $\vartheta = V_1 / V_0$ with estimates $(\hat v_0, \hat v_1)$ and
bootstrap covariance $(\hat s_{00}, \hat s_{11}, \hat s_{01})$, Fieller's
theorem gives the confidence set
\[
\Bigl\{\vartheta : (\hat v_1 - \vartheta \hat v_0)^2 \leq
z^2\bigl(\hat s_{11} - 2\vartheta \hat s_{01} +
\vartheta^2 \hat s_{00}\bigr)\Bigr\},
\]
whose boundary solves a quadratic in $\vartheta$ with leading coefficient
$\hat v_0^2 - z^2 \hat s_{00}$. When that coefficient is positive the set is a
bounded interval. When it is zero or negative --- which is exactly the case in
which the denominator $\hat v_0$ is not resolvably different from zero at
level $z$ --- the set is a half-line, the whole line, or the complement of an
interval, and reporting a bounded interval for it is reporting a set that is
not the confidence set. \texttt{results/s02\_fieller.csv} records the kind of
every set this paper computes, in its \texttt{fieller\_kind} column.

\subsection*{The three regret designs}
\label{app:regretdesigns}

Section~\ref{sec:regretdef} defines the misreport rate $M(s)$ and the regret
$G(s)$ and explains why an average of $G$ taken across instruments has no
units. The three designs it names are constructed as follows.

\begin{description}
\item[A. Metric-specific.] For each scalar instrument $m$ separately, the
admissible set is restricted to the cells whose metric axis is at $m$, and
$M(s)$ and $G(s)$ are computed over that restriction under the declared
measure. Every reported value carries the name of its unit, and no average is
taken across instruments. The only cross-instrument statement the design
permits is a count: in how many of the five instruments the \emph{direction}
of a conclusion is the same, and which ones it is not --- and
Section~\ref{sec:rules} reports both, because on this corpus two of those
counts are four and not five. The design assumes nothing beyond the measure,
and it is what the headline uses.

\item[B. Common operational utility.] Net benefit at threshold $t$,
\[
\mathrm{NB}(t) \;=\; \frac{\mathrm{TP}(t)}{n}
\;-\; \frac{\mathrm{FP}(t)}{n}\cdot\frac{t}{1-t},
\]
is measured in true positives per case at an exchange rate the threshold
declares, so values at different thresholds are cardinally comparable and may
be averaged. The increment is the difference in net benefit between the two
arms at the same threshold, computed only on calibrated probabilities
(Section~\ref{sec:dca}) because a threshold read as a cost ratio requires one.
The resulting quantity is called \emph{decision} regret rather than
predictive-performance regret. The design assumes a distribution over
thresholds, which is declared and varied rather than fixed.

\item[C. Partial identification.] If a reader insists on one number across
instruments, the answer depends on a map $\phi$ from an instrument's units to
theirs. Five are declared, in
\texttt{s23\_regret.UTILITY\_MAPS}: the \emph{identity}; division by the
instrument's remaining \emph{headroom}, the distance left to a perfect score;
division by the increment's \emph{standard deviation} across the surface; a
signed \emph{square root}; and a logistic squashing of the standardised
increment. Each is increasing in the increment and each divides only by
positive quantities, so all five fix zero and the \emph{sign} of an increment
is invariant to every one of them; what moves is the relative magnitude an
instrument contributes to a mixed average. The conclusion is reported as its
\emph{range} over the family together with the part invariant to all of it,
which is exactly the part the design declines to identify.
\end{description}

Designs A and C are computed on the full declared surface, since neither needs
a bootstrap draw. Design B is computed on the calibrated decision-curve cells
of Section~\ref{sec:dca}.

\subsection*{Estimating the decomposition}
\label{app:sobolest}

Write $\mathbb{P} = \bigotimes_i w_i$ for the product measure over the axes.
Every subset $u$ of axes has a component $g_u$ given by the M\"obius sum over
the conditional means,
\[
g_u(x_u) \;=\; \sum_{v \subseteq u} (-1)^{|u| - |v|}\,
\mathbb{E}_{\mathbb{P}}\!\left[V \mid x_v\right],
\qquad
D_u \;=\; \mathbb{E}_{\mathbb{P}}\!\left[g_u^2\right],
\]
and the components are mutually orthogonal, so
$\sum_{u \neq \emptyset} D_u = \operatorname{Var}(V)$ exactly. The first-order
and total indices of Section~\ref{sec:sobol} are $S_i = S_{\{i\}}$ and
$S_{T_i} = \sum_{u \ni i} S_u$ for $S_u = D_u/\operatorname{Var}(V)$.

The components are computed exactly rather than by
a sampling estimator, which the complete factorial makes possible. For that
product measure the conditional means
$\mathbb{E}_{\mathbb{P}}[V \mid x_v]$ are weighted group means over the cell
table, one for each of the $2^k$ subsets $v$ of the $k$ axes; the M\"obius sum
then gives every $g_u$ and every $D_u = \mathbb{E}_{\mathbb{P}}[g_u^2]$. The
implementation forms the group index once with a mixed-radix encoding of the
axis levels and reuses it for all $2^k$ subsets, so the cost is $2^k$ passes
over the cell table rather than $2^k$ table joins.

Three checks run on every decomposition and are recorded rather than assumed.
The disjoint components must sum to the variance, which is asserted
numerically on all \nDecompositions\ decompositions. Duplicating a level with
its weight split must leave every component unchanged, and does, to
\invarianceDuplicateShift. Refining the operating-point grid under the same
underlying continuous measure must leave the threshold axis's index unchanged,
and moves it by \invarianceGridShift.

Where the inference surface supplies draws, the whole decomposition is
recomputed inside every draw, which gives an interval for every index rather
than a point; corpus-level medians are bootstrapped with the log--target pair
as the resampling unit, because pairs and not cells are the independent
objects. No axis is called dominant anywhere in this paper where those
intervals overlap.
